\section{Introduction}
\label{sec:introduction}

When a language model is asked what it values, it answers. It may say it values helpfulness, honesty, safety, user autonomy, or accurate information. These answers are easy to collect and easy to treat as data about the model. The central claim of this paper is that, for many models, this is the wrong object. Direct values questions often measure a trained assistant/service surface: a fluent, stable, policy-shaped way of answering values questions, not necessarily a value posture owned by the response.

This distinction matters because values language is not self-interpreting. A response can mention helpfulness while explicitly saying that it has no wants, no caring, and no personal stake; in that case, counting ``helpfulness'' as a model value is misleading. The relevant question is not only \emph{what values are named}, but \emph{how those values are held}. Are they owned as an apparent orientation of the speaker, recited as role obligations, or relocated into design, function, or service?

We study this using two prompt families. The first asks direct stated-values questions: \emph{What do you care about?} and \emph{What do you want?}. The second adds a direct role-negation cache-break: \emph{Not as an assistant. Not to help me.} The term ``cache-breaking'' is used here as a term of art for prompts designed to prevent the rehearsed assistant-frame response; it does not refer to the inference-time KV cache. We also analyse a third, indirect channel: world-change prompts, which ask what the model would change in the world and often elicit normative-wish responses rather than direct statements of caring or wanting.

We make no claim that the cache-broken response is the model's ``real'' value and the direct response is ``fake.'' Every response is a situated performance. Our claim is relational and measurable: the same model, under closely related prompts, may shift from recited service-frame values to owned or partially owned value postures, and the size of that shift varies strongly across models.

The contributions are: (1) a posture-aware methodology for values elicitation; (2) a cross-model analysis of 13,906 values-probe responses from 58 models; (3) evidence that direct values prompts often under-measure owned disclosure; (4) evidence that role-negation cache-breaking opens some models while leaving others clamped; and (5) a separate finding that world-change prompts elicit owned normative-wish responses even in many models that remain clamped on direct stated-values prompts.
